\documentclass[lettersize,journal]{IEEEtran}
\IEEEoverridecommandlockouts
\usepackage{amsmath,amsfonts}
\usepackage{algorithmic}
\usepackage{algorithm}
\usepackage{array}
\usepackage[caption=false,font=normalsize,labelfont=sf,textfont=sf]{subfig}
\usepackage{textcomp}
\usepackage{url}
\usepackage{verbatim}
\usepackage{graphicx}
\usepackage{cite}
\usepackage{graphicx}
\usepackage{float}
\usepackage{booktabs}
\usepackage{picinpar}
\usepackage{amsmath}
\usepackage{url}
\usepackage[latin1]{inputenc}
\usepackage{colortbl}
\usepackage{soul}
\usepackage{multirow}
\usepackage{pifont}
\usepackage{color}
\usepackage{alltt}
\usepackage[hidelinks]{hyperref}
\usepackage{enumerate}
\usepackage{siunitx}
\usepackage{breakurl}
\usepackage{epstopdf}
\usepackage{pbox}
\usepackage{amssymb}
\usepackage{amsmath}
\usepackage{algorithm}
\usepackage[noend]{algpseudocode}
\usepackage{amsmath,amssymb,amsfonts}
\usepackage{graphicx}
\usepackage{booktabs}
\usepackage{multirow}
\usepackage{xcolor}
\usepackage{url}
\usepackage{etoolbox}

% Keep biography spacing fixed; the default 1fil glue spreads short columns.
% Patch the optional-argument implementation, not the public wrapper.
\makeatletter
\AtBeginDocument{%
  \def\@IEEEBIOphotowidth{0.78in}%
  \def\@IEEEBIOphotodepth{0.98in}%
  \def\@IEEEBIOhangwidth{0.90in}%
  \def\@IEEEBIOhangdepth{0.98in}%
  \def\@IEEEBIOskipN{6pt}%
  \expandafter\patchcmd\csname\string\IEEEbiography\endcsname
    {\vskip \@IEEEBIOskipN plus 1fil minus 0\baselineskip}
    {\vskip \@IEEEBIOskipN\relax}
    {\typeout{IEEE biography spacing compacted.}}
    {\PackageError{real_artical92}{Unable to compact IEEE biography spacing}
      {Check the IEEEtran biography definition before submission.}}%
}
\makeatother

\title{SafeAnchor-Grasp: An Object Grasp Optimization Method for Cluttered Multi-Object Scenes}

\author{Jiarui Hu, Xingyu Li, Rong Jiang, Zhipeng Wang,~\IEEEmembership{Member, IEEE},
Yanmin Zhou,~\IEEEmembership{Member, IEEE}, and Bin He,~\IEEEmembership{Senior Member, IEEE}%
\thanks{This work was supported in part by the National Key Research and Development Program of China under Grant 2024YFB4709800, in part by the National Natural Science Foundation of China under Grants 62688201, 62473294, 62403363, and 62403366, and in part by the China Postdoctoral Science Foundation under Grants 2024M762407 and 2026T190469.}%
\thanks{Corresponding author: Rong Jiang.}%
\thanks{Jiarui Hu, Xingyu Li, Rong Jiang, Zhipeng Wang, Yanmin Zhou, and Bin He are with the College of Electronics and Information Engineering, Tongji University, Shanghai 201804, China, and also with the National Key Laboratory of Autonomous Intelligent Unmanned Systems (Tongji University), Shanghai 201210, China, and also with the Frontiers Science Center for Intelligent Autonomous Systems, Shanghai 201210, China (e-mails: 2310203@tongji.edu.cn; xy\_l@tongji.edu.cn; 1810357@tongji.edu.cn; wangzhipeng@tongji.edu.cn; yanmin.zhou@tongji.edu.cn; hebin@tongji.edu.cn).}%
}

% \author{Anonymous Authors%
% \thanks{This work is prepared in the style of IEEE Robotics and Automation Letters (RA-L).}%
% }

\begin{document}
\maketitle

\begin{abstract}
In cluttered multi-object scenes, grasping requires not only a geometrically feasible hand pose but also stable physical interaction during closure, lifting, and holding. Existing methods can predict reasonable initial grasps, yet execution may still fail because of insufficient local contact or unstable closure. We therefore propose SafeAnchor-Grasp, a physical-executability enhancement framework for grasp-pose generation methods in cluttered multi-object scenes, with the goal of improving execution from target contact to stable lifting. The framework contains two core modules, AnchorLift and Rescue-Risk Gate (R\textsuperscript{2}-Gate). To improve the contact quality and lifting stability of execution-weak grasps, AnchorLift combines a Raw Grasp with the target-local point cloud, constructs Structured Physical Supervision (SPS), and generates same-target candidates. R\textsuperscript{2}-Gate jointly models the rescue benefit and degradation risk of each candidate to determine whether it should be adopted. Over 90 test scenes, SafeAnchor-Grasp improves the grasp success rates of DexDiffuser, GenDexGrasp, and DexGraspNet 2.0 by approximately 15, 10, and 3 percentage points, respectively, validating its ability to enhance physical executability across grasp generators.
\end{abstract}

\begin{IEEEkeywords}
Dexterous grasping, cluttered scenes, same-target enhancement, selective intervention, structured physical supervision.
\end{IEEEkeywords}
\begin{figure}[t]
  \centering
  \includegraphics[width=\linewidth]{figure1}
  \caption{Physical grasp optimization objective of SafeAnchor-Grasp in cluttered multi-object scenes. Given a scene observation, a raw grasp may already be close to the specified target yet still fail to form a stable closure because of insufficient local contact. SafeAnchor-Grasp uses Structured Physical Supervision (SPS) to enhance grasp candidates while preserving the target object, producing richer structured contacts and a stable closure. R\textsuperscript{2}-Gate then evaluates the candidate and the raw grasp as a pair to decide whether to retain the original grasp or perform the replacement.}
  \label{fig:intro}
\end{figure}
\section{Introduction}

In industrial sorting and target-specific picking tasks, robots must accurately identify and stably manipulate target objects in cluttered multi-object scenes. Compared with conventional two-finger grippers, dexterous hands have higher degrees of freedom, can form richer contact patterns on object surfaces, and can adapt to the complex geometry of irregular objects. This capability not only determines whether a robot can stably lift an object but also affects subsequent fine manipulation~\cite{ref1},~\cite{ref2}.

However, dexterous grasping in cluttered multi-object scenes remains difficult. Incomplete observations and interference from neighboring objects restrict the accessible space of the palm and fingers and weaken the contact quality and lifting stability of grasp candidates. A geometrically plausible grasp may still fail during execution because of insufficient contact, unstable closure, or post-lift instability. In complex scenes, failure does not always arise from selecting a completely wrong target. More often, the grasp has already reached or contacted the target object but still fails to form stable physical interaction during closure, lifting, and holding. This indicates that a clear gap remains between a reasonable grasp pose and stable physical execution.

Recent dexterous grasp generation methods have substantially improved grasp-pose prediction. With point-cloud features, generative models, and large-scale simulation data, existing methods can produce strong initial grasps in complex scenes. Zhang \emph{et al.} combined 6-D pose estimation with protective correction to reduce grasp collisions caused by misrecognition~\cite{ref3}. Liu \emph{et al.} developed continual grasp detection through knowledge transfer to improve adaptation to newly encountered objects and scenes~\cite{ref4}. These initial grasps usually contain valuable target-selection, hand-pose, and local-geometry priors. They therefore should not be simply discarded; instead, these priors can support unified physical-executability enhancement so that grasps that have already reached or contacted the target can achieve stable closure, lifting, and holding. Optimizing around existing grasps preserves these priors and corrects execution failures during contact, closure, and lifting, thereby providing a unified means of enhancing physical executability for different grasp-pose generation methods. However, an optimized candidate is not necessarily better than the original grasp. A new candidate may rescue a failed Raw Grasp, but it may also disrupt a Raw Grasp that would otherwise execute successfully. In cluttered multi-object scenes, candidate selection may also deviate from the original target when target-consistency constraints are absent. Physical-executability enhancement must therefore not only improve contact quality and lifting potential on the same target object, but also determine when optimization produces an effective rescue while avoiding new degradation.

Based on the above analysis, we propose SafeAnchor-Grasp, a physical-executability enhancement framework for multiple grasp-pose generation methods in cluttered multi-object scenes. It addresses the physical-execution ``last mile'' from reaching or contacting a target to stable closure, lifting, and holding. As shown in Fig.~\ref{fig:intro}, SafeAnchor-Grasp repairs execution-weak grasps while preserving target identity. It uses an existing Raw Grasp as the optimization basis and combines same-target candidate enhancement with risk control. Experiments across multiple grasp generators show that SafeAnchor-Grasp effectively improves Raw Grasp success.


The main contributions of this paper are as follows:
\begin{enumerate}
    \item We develop AnchorLift as a candidate enhancement module. AnchorLift uses a Raw Grasp as an anchor seed and combines it with the target-object local point cloud to generate an AnchorLift candidate. We introduce Structured Physical Supervision (SPS), coupled with the dexterous hand--object configuration, to improve candidate contact quality, hand posture, and lifting potential while preserving the original target object and initial grasp prior, thereby raising the same-object candidate ceiling.
    \item We develop Rescue-Risk Gate (R\textsuperscript{2}-Gate), a degradation-aware selective intervention mechanism. R\textsuperscript{2}-Gate models the Raw Grasp and AnchorLift candidate as a pair, learns candidate-success, rescue, and degradation confidence scores from candidate-level structured metadata, and drives risk-aware replacement using the intervention-score difference between the candidate and Raw Grasp. This preserves valuable existing grasps while exploiting candidate complementarity.
    \item We validate physical-executability enhancement across multiple grasp generators. We evaluate SafeAnchor-Grasp with multiple grasp generators over 90 GraspNet-1Billion test scenes. The results show that the framework improves grasp success across multiple generators, validating its ability to enhance physical executability for different grasp-pose generation methods. Further ablation experiments and real-robot experiments verify the effectiveness of its constituent modules and its applicability in a real perception--decision--execution pipeline, respectively.
\end{enumerate}


\section{Related Work}

\subsection{Dexterous Grasp Generation in Cluttered Scenes}

Grasping research has long studied how to generate executable grasp candidates from point clouds. Early methods mainly focused on parallel-jaw grippers and planned grasps by detecting 6-DoF grasp poses, predicting grasp quality or evaluating local point clouds, including GPD~\cite{ref5}, Dex-Net~\cite{ref6}, GG-CNN~\cite{ref7}, 6-DoF GraspNet~\cite{ref8}, and PointNetGPD~\cite{ref9}. Large-scale datasets and benchmarks further advanced cluttered-scene grasping. Dex-Net 4.0~\cite{ref10}, GraspNet-1Billion~\cite{ref11}, Contact-GraspNet~\cite{ref12}, VGN~\cite{ref13}, and AnyGrasp~\cite{ref14} demonstrate that learning-based methods can generate stable candidates in complex scenes. MetaGraspNetV2 provides a bin-picking benchmark with object relationship reasoning and dexterous grasping annotations~\cite{ref15}, occlusion-aware grasping reasons about occlusion relations and multi-target grasp poses in complex scenes~\cite{ref16}, and RIGNet focuses on intended-target grasping in densely stacked scenes through RoI-based intention and grasp reasoning~\cite{ref17}. Dexterous hand grasping is more difficult than parallel-jaw grasping because the system must predict both the palm pose and a high-dimensional joint configuration while forming stable multi-finger contacts during closure and lifting. For dexterous grasping under partial observation, Zhao \emph{et al.} studied 6-DoF dexterous grasp prediction in cluttered scenes from single-view point clouds~\cite{ref18}. Recent work uses large-scale simulation data, contact representations, diffusion models, and transformer set prediction~\cite{ref19} to learn dexterous grasp distributions. DexGraspNet~\cite{ref20}, GenDexGrasp~\cite{ref21}, UniDexGrasp~\cite{ref22}, DexGraspNet 2.0~\cite{ref23}, DexDiffuser~\cite{ref24}, DexGrasp Anything~\cite{ref25}, G-DexGrasp~\cite{ref26}, and ClutterDexGrasp~\cite{ref27} extend dexterous grasp candidate spaces in complex scenes from different perspectives. These methods mainly aim to enlarge or improve the grasp candidate space. In contrast, this paper does not treat the generator output directly as the final grasp. Instead, it uses AnchorLift as a same-object enhancement module that generates AnchorLift candidates bound to raw grasp. Final performance depends on whether R\textsuperscript{2}-Gate can identify which AnchorLift candidates should replace raw grasp and which raw grasp samples should be retained.

\begin{figure*}[t]
  \centering
  \includegraphics[width=\textwidth]{figure2}
  \caption{Complete pipeline of SafeAnchor-Grasp. (a) The data construction stage matches the raw grasp to the original target object and builds same-object enhancement inputs. (b) AnchorLift uses the target point cloud, local point cloud, raw grasp pose and qpos to generate an AnchorLift candidate. (c) R\textsuperscript{2}-Gate compares raw grasp and candidate features, predicts success, rescue potential and degradation risk and makes a keep-or-replace decision.}
  \label{fig:pipeline}
\end{figure*}

\subsection{Grasp Evaluation}

A geometrically reasonable grasp candidate does not necessarily succeed during execution. Therefore, many studies further investigate grasp quality evaluation, candidate refinement and test-time optimization. For parallel-jaw grasping, Dex-Net 2.0~\cite{ref6} uses GQ-CNN to predict grasp quality, PointNetGPD~\cite{ref9} evaluates local gripper point clouds, and 6-DoF GraspNet~\cite{ref8} combines grasp sampling with evaluator refinement by filtering and adjusting samples with a learned evaluator. For human or dexterous hand grasping, DexDiffuser~\cite{ref24}, ContactOpt~\cite{ref28}, GraspTTA~\cite{ref29}, and ContactGen~\cite{ref30} improve candidate quality from contact consistency, hand-object relation optimization, generative refinement, or evaluator-guided sampling. CADGrasp~\cite{ref31} further uses contact- and collision-aware intermediate representations to improve contact and collision avoidance for dexterous grasping in cluttered scenes. These studies show that a grasping system still needs additional evaluation or refinement after candidate generation. However, existing methods usually ask whether a single candidate itself is better, for example whether it has a higher quality score, a more reasonable contact region, lower collision risk, or stronger hand-object geometric consistency. For a system that already has a strong raw grasp, judging only the absolute quality of a new candidate is insufficient. A new candidate may improve contact and posture in failure cases and therefore rescue the original failure, but it may also introduce pose deviation, unstable closure, or local contact degradation in cases where raw grasp would have succeeded. In addition, candidate selection in cluttered multi-object scenes faces an object-consistency issue. If arbitrary candidate-pool selection is allowed, a higher oracle success may come from switching to an easier object rather than truly improving the grasp on the original target object.

From the perspective of the decision variable, existing candidate evaluation mainly learns an absolute quality function $Q(g^{a})$, whereas this work focuses on the replacement value $V(g^{0},g^{a})$ relative to an existing output. The latter must consider both whether a candidate converts a raw failure into success and whether it converts a raw success into failure. The objective is therefore not to further enlarge the candidate pool, but to learn when an enhanced candidate should replace the existing grasp.

\section{Method}

Given an initial visual grasp in a cluttered multi-object scene, SafeAnchor-Grasp neither discards the original output nor reselects an arbitrary successful grasp from a candidate pool. Instead, it treats raw grasp as a reusable anchor prior. The system first binds raw grasp to its original target object and then uses AnchorLift to generate an AnchorLift candidate on the same object. Because an AnchorLift candidate is not necessarily safer than raw grasp, R\textsuperscript{2}-Gate evaluates the two as a pair. The risk-aware replacement rule retains raw grasp by default and replaces it only when the candidate satisfies same-object and low-degradation-risk constraints.

\subsection{Same-Object Anchor Binding and Dataset Construction}

Same-object anchor binding is the boundary condition of SafeAnchor-Grasp. As shown in Fig.~\ref{fig:pipeline}(a), we first bind each raw grasp to its original target object and complete multi-view scene point-cloud matching to reduce viewpoint-dependent input differences in real scenes. This binding is then used as a hard constraint for subsequent candidate enhancement and risk-aware replacement. Each raw grasp anchor records scene id, view id, object id, raw grasp, confidence and available local point cloud.

For sample $i$, the same-object anchor is written as
\setcounter{equation}{0}
\begin{equation}
b_i = (s_i, v_i, o_i, c_i, g_i^{0}),
\quad
g_i^{0}=[t_i^{0}, r_i^{0}, q_i^{0}],
\label{eq:anchor_binding}
\end{equation}
where $s_i$ and $v_i$ denote scene and view indices, $o_i$ is the original target-object identity, $c_i$ is the object or crop center and $g_i^{0}$ is the raw grasp anchor.

The target-local point cloud is determined by object binding. This design ensures that AnchorLift inputs, training targets and final replacement all revolve around the object originally associated with raw grasp. We construct a scene-object linked scene-level dataset by recording the selected object, multi-view information, multi-ground-truth information, viewed-object point cloud, raw grasp and local point cloud. The key of this stage is not to expand arbitrary object candidates, but to establish a stable and reproducible input boundary for same-object enhancement.

The target-local crop used by AnchorLift is defined by the  raw grasp-bound object,
\begin{equation}
P_i^{o} =
\{p_k \in P_i \mid \mathrm{seg}(p_k)=o_i,\ 
\|p_k-c_i\|_2 \le r_c \},
\label{eq:target_crop}
\end{equation}
where $P_i$ is the scene point cloud, $\mathrm{seg}(p_k)$ denotes the object assignment of point $p_k$ and $r_c$ is the local crop radius. This definition prevents the candidate generation stage from silently switching to a different object.

\subsection{AnchorLift: Anchor-Conditioned Grasp Optimization }

AnchorLift aims to work together with raw grasp to raise the same-object candidate ceiling. As shown in Fig.~\ref{fig:pipeline}(b), it uses raw grasp as an anchor seed and adjusts hand pose, local contact location and joint configuration on the same target object, rather than generating an unconstrained grasp from scratch. Its inputs include the target-object point cloud, target-local point cloud, raw grasp pose, qpos and object binding information, and its output is an AnchorLift candidate. Importantly, AnchorLift should not be interpreted as a small local refinement of the raw grasp pose or as simple post-processing around raw grasp. Its goal is to expand the candidate space under the same-object anchor constraint. The AnchorLift mapping is written as
\begin{equation}
g_i^{a} = f_{\theta}(P_i^{o}, g_i^{0}),
\quad
g_i^{a}=[t_i^{a}, r_i^{a}, q_i^{a}],
\label{eq:anchorlift_generation}
\end{equation}
where $f_{\theta}$ generates an AnchorLift candidate from the target-local geometry and raw grasp seed.

AnchorLift uses a target-local point encoder, a seed-grasp encoder, and a diffusion candidate decoder. Its input features are organized into four branches: the target-object point cloud, target-local point cloud, raw-grasp end-effector pose, and hand qpos are encoded separately and then fused. Each training sample contains a 2,048-point target-local cloud cropped around the center selected by raw grasp and a 25-dimensional raw seed comprising 3D translation, a 6D rotation representation, and 16 dexterous-hand joint angles. A PointTransformer encodes the local cloud into a scene token. A two-layer multilayer perceptron (MLP) maps the 25-dimensional raw seed to a 512-dimensional seed token, which conditions the cross-attention diffusion network together with the scene token. During training, SPS comprises the Hand-Structure Field (HSF), Surface-Support Field (SSF), and Penetration Field (PF) and is optimized jointly with the direct pose-regression loss. HSF promotes a valid hand-skeleton structure, SSF encourages local surface support, and PF suppresses hand--object penetration.

\begin{figure}[t]
  \centering
  \includegraphics[width=\linewidth]{figure3}
  \caption{Core intermediate states of AnchorLift and R\textsuperscript{2}-Gate. The left side shows raw grasp and the right side shows the same-target AnchorLift candidate selected by R\textsuperscript{2}-Gate. The colored overlays diagnose the hand structure associated with HSF, target-surface proximity associated with SSF, and hand--object penetration associated with PF, providing a geometric interpretation of the three SPS constraints.}
  \label{fig:diagnostic}
\end{figure}

Fig.~\ref{fig:diagnostic} compares raw grasp and the same-target AnchorLift candidate selected by R\textsuperscript{2}-Gate. The colored overlays diagnose the hand structure associated with HSF, target-surface proximity associated with SSF, and hand--object penetration associated with PF, thereby providing a geometric interpretation of the three SPS constraints.

For sample $i$, let $\mathbf{x}_{0,i}\in\mathbb{R}^{25}$ denote a successful grasp matched to the raw seed under the same scene, view, and object conditions, and let $\boldsymbol\epsilon\sim\mathcal N(\mathbf 0,\mathbf I)$ be diffusion noise. At a randomly sampled time step $t$, the forward process constructs
\begin{equation}
\mathbf{x}_{t,i}=\sqrt{\bar\alpha_t}\,\mathbf{x}_{0,i}
+\sqrt{1-\bar\alpha_t}\,\boldsymbol\epsilon .
\tag{4a}\label{eq:diffusion_forward}
\end{equation}
The network predicts noise under the point-cloud and raw-seed condition $\mathbf c_i$. The current configuration uses the L1 noise objective
\begin{equation}
\mathcal L_{\mathrm{DFM}}
=\mathbb E_{i,t,\boldsymbol\epsilon}
\left[\left\|\boldsymbol\epsilon-
\boldsymbol\epsilon_\theta(\mathbf{x}_{t,i},t,\mathbf c_i)\right\|_1\right].
\tag{4b}\label{eq:dfm_loss}
\end{equation}

Rather than constraining only the 25-dimensional parameter space, SPS passes the recovered denoised grasp $\hat{\mathbf{x}}_{0,i}$ through differentiable forward kinematics. This yields hand-surface points $\{\mathbf h_j\}$, contact-distance keypoints $\{\mathbf d_k\}$, and hand-structure keypoints $\{\mathbf u_m\}$. The target-local point cloud supplies object points $\mathbf p$ and normals $\mathbf n$. HSF, SSF, and PF are constructed online from these deterministic geometric quantities and require no additional manually annotated field labels.

PF describes the penetration constraint between the hand surface and the target object. For each hand-surface point $\mathbf h_j$, let $\mathbf p_j^*$ be its nearest target-cloud point and $\mathbf n_j^*$ the corresponding normal. PF uses the normal sign to identify hand points on the interior side of the object and penalizes the maximum penetration depth in each sample:
\begin{equation}
\mathcal L_{\mathrm{PF}}
=\frac{1}{B}\sum_{i=1}^{B}
\max_j\left[
\mathbb I\!\left((\mathbf p_{ij}^*-\mathbf h_{ij})^\top
\mathbf n_{ij}^*>0\right)
\left\|\mathbf p_{ij}^*-\mathbf h_{ij}\right\|_2
\right].
\tag{4c}\label{eq:pf_loss}
\end{equation}

SSF describes local support near the target surface and encourages contact-distance keypoints to approach that surface. Let $r_{ik}=\min_{\mathbf p\in P_i^o}\|\mathbf d_{ik}-\mathbf p\|_2$ and $\mathcal A=\{(i,k):r_{ik}<r_c\}$, where the current implementation uses $r_c=0.02\,\mathrm m$:
\begin{equation}
\mathcal L_{\mathrm{SSF}}
=\frac{\sum_{(i,k)\in\mathcal A}r_{ik}}
{|\mathcal A|+10^{-5}}.
\tag{4d}\label{eq:ssf_loss}
\end{equation}

HSF constrains the internal structure of the dexterous hand and suppresses distances below a safety margin between distinct hand keypoints. The current implementation uses $r_s=0.02\,\mathrm m$:
\begin{equation}
\mathcal L_{\mathrm{HSF}}
=\frac{1}{B}\sum_{i=1}^{B}\sum_{m\ne n}
\left[r_s-\|\mathbf u_{im}-\mathbf u_{in}\|_2\right]_+ .
\tag{4e}\label{eq:hsf_loss}
\end{equation}

SPS therefore comprises three explicit physical constraints, $\mathcal L_{\mathrm{SPS}}=\mathcal L_{\mathrm{HSF}}+\mathcal L_{\mathrm{SSF}}+\mathcal L_{\mathrm{PF}}$. Together with the diffusion noise-prediction loss $\mathcal L_{\mathrm{DFM}}$ and the direct pose-regression loss $\mathcal L_{\mathrm{POS}}$, the total AnchorLift loss is:
\begin{equation}
\begin{aligned}
\mathcal L_{\mathrm{Anchor}}
&=\mathcal L_{\mathrm{DFM}}+\mathcal L_{\mathrm{POS}}
+\mathcal L_{\mathrm{HSF}}\\
&\quad+\mathcal L_{\mathrm{SSF}}+\mathcal L_{\mathrm{PF}}\\
&=\mathcal L_{\mathrm{DFM}}+\mathcal L_{\mathrm{POS}}
+\mathcal L_{\mathrm{SPS}}.
\end{aligned}
\tag{4f}\label{eq:anchorlift_loss}
\end{equation}

\subsection{R\textsuperscript{2}-Gate: Risk-Aware Grasp Replacement Selection}

Rescue-Risk Gate (R\textsuperscript{2}-Gate) evaluates whether adopting a same-target AnchorLift candidate is preferable to retaining raw grasp. For the $i$th scene--view group, let $g_i^0$ denote raw grasp and $g_i^a$ the AnchorLift candidate.

R\textsuperscript{2}-Gate does not directly encode the local point cloud. Instead, it constructs 82-dimensional candidate-level metadata for raw grasp and the AnchorLift candidate, each referenced to raw grasp. The input summarizes candidate and raw scores and ranks, candidate source, relative pose and joint changes, target identity, and geometric availability. CandidateNet uses an $82\!\rightarrow\!128\!\rightarrow\!128$ metadata MLP and a $128\!\rightarrow\!256\!\rightarrow\!128$ fusion MLP, with 93,443 trainable parameters in total.

CandidateNet outputs candidate-success, rescue, and degradation confidence scores. Let $y_i^a,y_i^0\in\{0,1\}$ denote the physical success labels of the AnchorLift candidate and raw grasp, respectively. The three supervision labels are
\begin{equation}
y_i^{\mathrm{succ}}=y_i^a,\qquad
y_i^{\mathrm{res}}=(1-y_i^0)y_i^a,\qquad
y_i^{\mathrm{deg}}=y_i^0(1-y_i^a).
\tag{5}\label{eq:r2_labels}
\end{equation}
The training objective contains success, rescue, and degradation classification losses together with ranking and raw-preserve terms,
\begin{equation}
\begin{aligned}
\mathcal L_{\mathrm{R2}}
=&\;\mathcal L_{\mathrm{succ}}^{\mathrm{BCE}}
+2\mathcal L_{\mathrm{res}}^{\mathrm{BCE}}
+5\mathcal L_{\mathrm{deg}}^{\mathrm{BCE}}\\
&+\frac{1}{2}\left(\mathcal L_{\mathrm{pair}}+\mathcal L_{\mathrm{list}}\right)
+\mathcal L_{\mathrm{pres}}.
\end{aligned}
\tag{6}\label{eq:r2gate_loss}
\end{equation}
The rescue and degradation terms support rescue identification and degradation suppression, respectively; ranking distinguishes successful from failed entries in the same group, and raw-preserve protects raw grasps that already succeed.

R\textsuperscript{2}-Gate computes an intervention score for raw grasp and the AnchorLift candidate and uses their difference for gating:
\begin{equation}
J_i(g)=s_i^{\mathrm{succ}}(g)+\alpha s_i^{\mathrm{res}}(g)-\beta s_i^{\mathrm{deg}}(g),
\qquad
m_i=J_i(g_i^a)-J_i(g_i^0).
\tag{7}\label{eq:intervention_margin}
\end{equation}
A candidate enters the eligible set only if it preserves the target, satisfies $m_i>\tau$, and has $s_i^{\mathrm{deg}}\leq\delta$. Eligible candidates are ranked in descending order of intervention-score difference for risk-aware replacement selection, and the selected candidates form the replacement set $R$. The final output is
\begin{equation}
g_i^{\mathrm{out}}
=
\begin{cases}
g_i^a,& i\in R,\\
g_i^0,& i\notin R.
\end{cases}
\tag{8}\label{eq:final_decision}
\end{equation}
\section{Experiments}

\subsection{Experimental Setup}
Experiments use the GraspNet-1Billion dataset~\cite{ref11}. AnchorLift is trained on scenes 0000--0099. R\textsuperscript{2}-Gate is trained using 23,028 same-target raw--AnchorLift pairs from scenes 0000--0089 and validated using 2,552 pairs from scenes 0090--0099. Scenes 0100--0189 are used exclusively for testing. They contain 90 scenes with 256 views per scene, giving 23,040 scene--view samples in total. Each scene--view produces one top-1 grasp and is executed in Isaac Gym~\cite{ref32}. A grasp is counted as successful when it passes the pre-grasp collision check and the bound target object is lifted by more than 0.10 m.

To evaluate both module compatibility and the performance of the module itself, the experiments are organized as follows. Section IV-B verifies SafeAnchor-Grasp as an insertable optimization module for different grasp generators. Section IV-C analyzes candidate enhancement, failure rescue, degradation control, and ablation results based on the baseline. Section IV-D deploys the grasping model on a real robot to verify its practical applicability.

\subsection{Lightweight Integration Test Across Grasp Generators}

We use DexDiffuser~\cite{ref24}, GenDexGrasp~\cite{ref21}, and DexGraspNet 2.0~\cite{ref23}, three heterogeneous dexterous-hand grasp generators, to evaluate the cross-source integration capability of SafeAnchor-Grasp. SafeAnchor-Grasp does not rewrite their network structures or parameters. Instead, their outputs are mapped to target-bound Raw Grasps, after which SafeAnchor-Grasp is inserted. Adaptation requires only lightweight source calibration and does not involve joint retraining of the generator or refinement backbone. The adaptation takes 42.87 s; therefore, from the perspective of model modification and deployment integration, the additional cost is negligible. Each Raw Grasp remains bound to its original target object and test view, and same-target candidates are generated around this Raw anchor. Switching target objects to obtain spurious performance gains is not allowed. This experiment evaluates whether SafeAnchor-Grasp can serve as a generator-decoupled post-processing enhancement module and provide plug-and-play integration while preserving different model interfaces and target priors.

\begin{table*}[t]
\centering
\caption{Performance comparison after integrating SafeAnchor-Grasp with different grasp generators. The table reports Raw Grasp success, final success after integration, absolute improvement, candidate replacement rate, and the joint grasp ceiling of Raw Grasp and the optimized candidates to assess cross-generator adaptation. Improvement is calculated from the reported success rates. Dashes denote joint-ceiling results pending recalculation.}
\label{tab:generator_adaptation}
\begingroup
\setlength{\tabcolsep}{8pt}
\renewcommand{\arraystretch}{1.35}
\small
\begin{tabular*}{\textwidth}{@{\extracolsep{\fill}}l c c c c c@{}}
\toprule
Grasp generator & Raw success & Success after integration & Improvement & Replacement rate & Raw--optimized joint ceiling \\
\midrule
DexDiffuser~\cite{ref24} & 0.73 & 0.88 & 0.15 & 98\% & -- \\
GenDexGrasp~\cite{ref21} & 0.79 & 0.89 & 0.10 & 93\% & -- \\
DexGraspNet 2.0~\cite{ref23} & 0.89 & 0.92 & 0.03 & 12\% & 0.95 \\
\bottomrule
\end{tabular*}
\endgroup
\end{table*}

Table~\ref{tab:generator_adaptation} shows that SafeAnchor-Grasp obtains positive gains for three grasp generators with different initial performance levels. Across 90 test scenes, DexDiffuser improves from 0.73 to 0.88, GenDexGrasp improves from 0.79 to 0.89, and DexGraspNet 2.0 improves from 0.89 to 0.92. At the precision reported in the table, the corresponding absolute gains are 15, 10, and 3 percentage points, respectively. All three generators can be connected to the unified enhancement framework without modifying their main bodies, indicating that the gain does not depend on the internal structure of a single generator.

The replacement-rate column of Table~\ref{tab:generator_adaptation} further characterizes the output selection of SafeAnchor-Grasp. The reported replacement rates of DexDiffuser and GenDexGrasp are 98\% and 93\%, respectively. Together with the success improvements, these results demonstrate cross-source adaptation and enhancement. These rates are specific to the respective generators and evaluation settings; the DexGraspNet 2.0 replacement statistics are analyzed separately in Sec.~\ref{sec:graspnet_analysis}. For DexGraspNet 2.0, the same-target Raw-or-AnchorLift ceiling on the test set is 94.81\%, consistent with Panel~B of Table~\ref{tab:core_ablation}; Table~\ref{tab:generator_adaptation} reports this value as 0.95 after rounding the decimal proportion to two decimal places. Its ceiling remains above the final selected success rate, indicating remaining candidate-selection potential for this generator.

\subsection{SafeAnchor-Grasp Analysis}
\label{sec:graspnet_analysis}
In the following experiments, we use DexGraspNet 2.0~\cite{ref23} as the baseline and evaluate SafeAnchor-Grasp on GraspNet-1Billion. All replacement rates reported in this subsection refer specifically to this generator and evaluation setting, rather than a shared replacement limit across generators.

\begin{figure*}[t]
  \centering
  \includegraphics[width=\textwidth]{figure4}
  \caption{Verified replay examples from low-success scenes. Each row shows the target object, a failed Raw Grasp, and the SafeAnchor-Grasp result on the same target object. The Raw Grasp usually reaches the vicinity of the target but cannot produce an effective lift, whereas the optimized same-object output achieves stable lifting. Red denotes the failed Raw Grasp and green denotes the selected SafeAnchor-Grasp output.}
  \label{fig:low_success_failure}
\end{figure*}

\subsubsection{SafeAnchor-Grasp Performance Test}

\begin{table}[t]
\centering
\caption{Overall and low-success scene performance. Low-success scenes are defined by raw grasp success below 0.85. SafeAnchor-Grasp achieves larger gains in these execution-weak scenes while preserving same-object output selection.}
\label{tab:overall_low_success}
\begingroup
\setlength{\tabcolsep}{1.8pt}
\renewcommand{\arraystretch}{1.35}
\scriptsize
\begin{tabular*}{\columnwidth}{@{\extracolsep{\fill}}l c c c c@{}}
\toprule
\textbf{Set} & \textbf{Scenes} & \textbf{Raw} & \textbf{SafeAnchor} & \textbf{$\Delta$} \\
\midrule
GraspNet-1Billion & 90 & 0.8921 & \textbf{0.9221} & \textbf{+0.0301} \\
Low-success scenes & 26 & 0.7237 & \textbf{0.8274} & \textbf{+0.1037} \\
Best-10 low-success & 10 & 0.6566 & \textbf{0.9008} & \textbf{+0.2441} \\
Best low-success cases & 6 & 0.6113 & \textbf{0.9277} & \textbf{+0.3164} \\
\bottomrule
\end{tabular*}
\endgroup
\end{table}

Table~\ref{tab:overall_low_success} first reports the overall result on 90 test scenes: SafeAnchor-Grasp improves mean success from 0.8921 to 0.9221, with a scene-level paired-bootstrap 95\% confidence interval of $[0.0138,0.0492]$. When evaluation is restricted to the 26 scenes whose raw-grasp success is below 0.85, mean success increases from 0.7237 to 0.8274, a gain of 10.37 percentage points, with a 95\% confidence interval of $[0.0550,0.1587]$. This difference indicates that SafeAnchor-Grasp is especially effective when the initial grasp prior preserves the target but remains physically weak during execution.

The best-10 and best-case rows are selected by gain from the same 26 low-success scenes rather than from additional test sets. They are used to examine the gains in execution-weak scenes most in need of repair. Representative scene-level success rates for scenes 0117, 0114, 0115, and 0121 improve from 0.4727, 0.6328, 0.6563, and 0.5195 to 0.9570, 0.9922, 0.9766, and 0.7734, respectively.

Fig.~\ref{fig:low_success_failure} shows verified replay examples from low-success scenes. Each row presents the target object, a failed Raw Grasp, and the SafeAnchor-Grasp result on the same target object. Raw Grasp is often already close to the target but fails because of insufficient lifting support or push/slide behavior during closure and lifting; SafeAnchor-Grasp repairs these physical execution failures while preserving the original target object.

\begin{table}[t]
\centering
\caption{Overall performance on GraspNet-1Billion. SafeAnchor-Grasp achieves the best performance under different evaluation settings.}
\label{tab:graspnet_benchmark}
\begingroup
\setlength{\tabcolsep}{2.6pt}
\renewcommand{\arraystretch}{1.5}
\small
\begin{tabular*}{\columnwidth}{@{\extracolsep{\fill}}c c c c c@{}}
\toprule
\multirow{2}{*}{Method} & \multirow{2}{*}{Ratio} & \multicolumn{3}{c}{GraspNet-1Billion} \\
\cmidrule(lr){3-5}
& & Dense & Random & Loose \\
\midrule
HGC-Net~\cite{ref33} & 1 & 46.0 & 37.8 & 26.7 \\
GraspTTA~\cite{ref29} & 1 & 62.5 & 54.1 & 42.8 \\
ISAGrasp~\cite{ref34} & 1 & 63.4 & 60.7 & 51.4 \\
DexGraspNet 2.0~\cite{ref23} & 1/1000 & 89.2 & 84.5 & 81.9 \\
CADGrasp~\cite{ref31} & 1/1000 & 89.5 & 87.2 & 81.6 \\
\textbf{SafeAnchor-Grasp} & \textbf{1/1000} & \textbf{92.2} & \textbf{90.1} & \textbf{85.6} \\
\bottomrule
\end{tabular*}
\endgroup
\end{table}

Table~\ref{tab:graspnet_benchmark} compares the overall performance of different methods on GraspNet-1Billion. Under the same 1/1000 ratio setting, SafeAnchor-Grasp reaches 92.2, 90.1, and 85.6 under the Dense, Random, and Loose settings, respectively. Compared with DexGraspNet 2.0~\cite{ref23}, the corresponding improvements are 3.0, 5.6, and 3.7 percentage points. Compared with CADGrasp~\cite{ref31}, SafeAnchor-Grasp improves by 2.7, 2.9, and 4.0 percentage points, respectively. This shows that same-object safe enhancement improves not only average raw-grasp performance but also stability under more challenging evaluation settings.

\subsubsection{AnchorLift Candidate Enhancement and R\textsuperscript{2}-Gate Risk-Aware Selection}

Fig.~\ref{fig:anchorlift} visualizes the selected object, target-local point cloud, and geometric diagnostics after SPS is applied, showing that AnchorLift consistently operates around the raw-grasp-bound target.

\begin{figure}[t]
  \centering
  \includegraphics[width=\linewidth]{figure5}
  \caption{AnchorLift intermediate states. The figure shows same-object binding and SPS-based geometric diagnostics for AnchorLift across different objects in the same scene.}
  \label{fig:anchorlift}
\end{figure}

Same-object candidate-ceiling statistics show that AnchorLift expands the candidate space while preserving target binding, but also produces both rescue and degradation cases. Thus, AnchorLift expands the candidate space, whereas R\textsuperscript{2}-Gate decides which candidates are worth adopting.

AnchorLift reaches same-object ceilings of 0.9819 on the training split and 0.9835 on the validation split. R\textsuperscript{2}-Gate ranks candidates by the intervention-score difference relative to Raw Grasp and selects outputs under target-consistency and degradation-risk constraints. On the DexGraspNet 2.0 test set, it replaces 2,764 of 23,040 outputs, yielding an observed replacement rate of approximately 12\%, while preserving the target object for every replacement. The final success rate increases from 89.21\% to 92.21\%.

The above DexGraspNet 2.0 test results show that R\textsuperscript{2}-Gate obtains a positive net gain within a limited intervention range. Replacement decisions must consider both rescue benefit and degradation risk rather than relying solely on a candidate's absolute success score.

\begin{figure}[t]
  \centering
  \includegraphics[width=\linewidth]{figure7}
  \caption{Real robotic arm and dexterous-hand grasping platform. The figure shows the real UR arm-dexterous hand grasping experimental platform.}
  \label{fig:scene}
\end{figure}

\begin{figure*}[t]
  \centering
  \includegraphics[width=\textwidth]{figure8}
  \caption{Real-robot execution demonstration of SafeAnchor-Grasp. The upper block shows observation, object feature, raw grasp, SafeAnchor-Grasp output and UR arm plus dexterous-hand physical execution for a single target object. The lower block shows the same workflow in a cluttered multi-object scene.}
  \label{fig:real_robot}
\end{figure*}

\subsubsection{Ablation Experiments}
\label{sec:ablation}

To separately verify candidate optimization and safe selection, Table~\ref{tab:core_ablation} divides the ablation study into two parts. Part~A examines whether AnchorLift depends on SPS, the target-local point cloud, and raw-grasp conditioning. Part~B compares never intervening, always intervening, Oracle intervention, and learned limited intervention with R\textsuperscript{2}-Gate.

\begin{table}[t]
\centering
\caption{Core module ablations of SafeAnchor-Grasp.}
\label{tab:core_ablation}
\begingroup
\setlength{\tabcolsep}{1.4pt}
\renewcommand{\arraystretch}{1.4}
\scriptsize
\begin{tabular*}{\columnwidth}{@{\extracolsep{\fill}}>{\centering\arraybackslash}p{0.42\columnwidth} c c c@{}}
\toprule
\multicolumn{4}{c}{\textbf{A: AnchorLift component ablation}} \\
\midrule
Variant & \shortstack{Diagnostic-scene\\success} & \multicolumn{2}{c}{\shortstack{Change\\vs. full}} \\
\midrule
\textbf{Full AnchorLift } & \textbf{57.00\%} & \multicolumn{2}{c}{--} \\
w/o SPS & 46.10\% & \multicolumn{2}{c}{$-10.90$\%} \\
w/o target-local point cloud & 42.70\% & \multicolumn{2}{c}{$-14.30$\%} \\
w/o raw-grasp condition & 39.40\% & \multicolumn{2}{c}{$-17.60$\%} \\
\midrule
\multicolumn{4}{c}{\textbf{B: R\textsuperscript{2}-Gate intervention-policy ablation}} \\
\midrule
Intervention policy & \shortstack{Test\\success} & \shortstack{Hard-scene\\success} & \shortstack{Replacement\\rate} \\
\midrule
Raw Grasp (never intervene) & 89.21\% & 72.37\% & 0\% \\
AnchorLift direct replacement (always intervene) & 84.31\% & 64.02\% & 100\% \\
Same-target Raw-or-AnchorLift ceiling & \textbf{94.81\%} & -- & Oracle \\
\textbf{SafeAnchor-Grasp (learned limited intervention)} & \textbf{92.21\%} & \textbf{82.74\%} & \textbf{12.00\%} \\
\bottomrule
\end{tabular*}
\endgroup
\end{table}

Panel~A of Table~\ref{tab:core_ablation} shows that SPS increases AnchorLift success on the diagnostic scenes from 46.10\% to 57.00\%. Removing the target-local point cloud or raw-grasp condition reduces success to 42.70\% and 39.40\%, respectively. These results show that SPS, target-local geometry, and the raw-grasp anchor jointly support same-target candidate optimization.

Panel~B of Table~\ref{tab:core_ablation} shows that unconditionally adopting AnchorLift reduces test success from 89.21\% to 84.31\%, confirming that the candidate set contains genuine degradation cases. Meanwhile, the same-target Oracle ceiling reaches 94.81\%, demonstrating substantial rescue potential in the same candidate set. In this DexGraspNet 2.0 setting, R\textsuperscript{2}-Gate achieves an observed replacement rate of approximately 12\%, reaches 92.21\% overall success, and raises hard-scene success to 82.74\%. Its role is therefore not to maximize the number of replacements, but to select reliable interventions from complementary candidates.

Further three-seed validation-set ablations show that removing the rescue component decreases rescue recall from 19.85\% to 13.33\% and net rescue from 12.41/1k to 7.58/1k. Removing the degradation component increases erroneous degradation from 5.09/1k to 8.75/1k and reduces selected-candidate success from 94.34\% to 89.43\%. These results show that the rescue and degradation components support effective-rescue identification and harmful-replacement suppression, respectively.

\subsection{Real-Robot Experiments}

Fig.~\ref{fig:scene} shows the real experimental platform, on which we quantitatively evaluate SafeAnchor-Grasp using a real UR robot arm and Allegro Hand. The system uses a RealSense D455 to acquire scene point clouds and RGB images and SAM~\cite{ref35} to segment candidate objects, thereby obtaining visible object point clouds and their surrounding local point clouds. After SafeAnchor-Grasp outputs the final grasp configuration, the system determines the target tool-center-point (TCP) pose, solves for the UR arm's joint configuration using inverse kinematics, and moves the Allegro Hand to the target grasp pose. This platform forms a complete robotic perception--computation--execution loop.

Fig.~\ref{fig:real_robot} shows the physical execution pipeline in single-object and cluttered multi-object scenes. Additional real-robot videos and project materials are available at \url{https://hujiarui2310203.github.io/SafeAnchor-Grasp/}.

The real-robot evaluation uses ten common everyday objects that do not appear in the training data. In the single-object setting, each of the same ten objects is grasped from ten randomized poses, producing 100 target-grasp attempts per method.

In the cluttered multi-object setting, all ten objects are randomly placed in the workspace at the beginning of each round. After grasping the specified target, the system removes it from the workspace and designates the next target from the remaining objects until each of the ten objects has undergone one grasp attempt. Each round therefore contains ten consecutive target-grasp attempts. After each round, all initial object poses are randomized again. The experiment is repeated for ten rounds, again producing 100 target-grasp attempts per method.

An attempt is counted as successful only when the specified target object is stably lifted by at least 0.1~m. This experiment demonstrates that SafeAnchor-Grasp outputs can be directly integrated into a real system while maintaining execution gains under object-pose variations and multi-object interference.

\begin{table}[t]
\centering
\caption{Quantitative real-robot comparison between Raw Grasp and SafeAnchor-Grasp. Each method performs 100 target-grasp attempts in both the single-object and cluttered multi-object settings. The table reports robot execution rate and final grasp success.}
\label{tab:real_robot_quantitative}
\begingroup
\setlength{\tabcolsep}{2.6pt}
\renewcommand{\arraystretch}{1.5}
\small
\begin{tabular*}{\columnwidth}{@{\extracolsep{\fill}}l c c c c@{}}
\toprule
\multirow{2}{*}{Setting} & \multicolumn{2}{c}{Raw Grasp} & \multicolumn{2}{c}{SafeAnchor-Grasp} \\
\cmidrule(lr){2-3}\cmidrule(lr){4-5}
 & \shortstack{Robot\\execution rate} & \shortstack{Grasp\\success rate} & \shortstack{Robot\\execution rate} & \shortstack{Grasp\\success rate} \\
\midrule
Single object & 0.83 & 0.72 & 0.92 & 0.88 \\
\shortstack[l]{Cluttered\\multi-object} & 0.67 & 0.52 & 0.77 & 0.69 \\
\bottomrule
\end{tabular*}
\endgroup
\end{table}

Table~\ref{tab:real_robot_quantitative} shows that SafeAnchor-Grasp improves both robot execution rate and final grasp success on the real platform. In the single-object setting, robot execution rate increases from 0.83 to 0.92 and grasp success from 0.72 to 0.88. In the cluttered multi-object setting, the two metrics increase from 0.67 and 0.52 to 0.77 and 0.69, respectively. The lower absolute performance in clutter reflects the additional effects of occlusion, neighboring-object interference, and sequential object removal, while SafeAnchor-Grasp maintains consistent gains in both settings.

\section{Conclusion}

This paper presented SafeAnchor-Grasp, which formulates dexterous grasp optimization in cluttered multi-object scenes as selective intervention under existing-grasp protection. Unlike direct grasp generation or unconditional adoption of new candidates, the method treats raw grasp as the current existing output, generates a same-target alternative action with AnchorLift, and uses R\textsuperscript{2}-Gate to jointly model positive rescue and negative degradation outcomes before deciding whether to retain raw grasp or adopt the optimized candidate.

With DexGraspNet 2.0 as the baseline, SafeAnchor-Grasp improves mean success across 90 GraspNet-1Billion test scenes from 0.8921 to 0.9221, a gain of 3.01 percentage points. On the 26 hard scenes whose raw success is below 0.85, mean success increases from 0.7237 to 0.8274. These results highlight the benefit of jointly evaluating the rescue benefit and degradation risk of each replacement rather than adopting optimized candidates indiscriminately. The system performs same-target intervention only when a candidate has high intervention value relative to retaining raw grasp, thereby converting candidate complementarity into a positive net gain.

Overall, SafeAnchor-Grasp reuses existing visual grasping priors without retraining the original grasp generator and uses risk-aware selection to determine when to adopt optimized candidates. It improves dexterous grasping in cluttered multi-object scenes without switching target objects or unconditionally overwriting raw grasp.

Quantitative real-robot experiments further show that SafeAnchor-Grasp outputs can be integrated into a UR arm--dexterous-hand pipeline and improve both robot execution rate and grasp success in single-object and cluttered multi-object settings. This modular design reduces the software and hardware modifications required for existing systems and provides an integrable technical route for flexible sorting, mixed-object handling, and reliable target grasping in cluttered environments.


\begin{thebibliography}{35}

\bibitem{ref1} D. Wang, C. Liu, F. Chang, N. Li and G. Li, ``High-Performance Pixel-Level Grasp Detection Based on Adaptive Grasping and Grasp-Aware Network,'' \emph{IEEE Transactions on Industrial Electronics}, vol. 69, no. 11, pp. 11611--11621, 2022.
\bibitem{ref2} S. Yu, D.-H. Zhai and Y. Xia, ``EGNet: Efficient Robotic Grasp Detection Network,'' \emph{IEEE Transactions on Industrial Electronics}, vol. 70, no. 4, pp. 4058--4067, 2023.
\bibitem{ref3} H. Zhang \emph{et al.}, ``A Practical Robotic Grasping Method by Using 6-D Pose Estimation With Protective Correction,'' \emph{IEEE Transactions on Industrial Electronics}, vol. 69, no. 4, pp. 3876--3886, 2022.
\bibitem{ref4} J. Liu, J. Xie, S. Huang, C. Wang and F. Zhou, ``Continual Learning for Robotic Grasping Detection With Knowledge Transferring,'' \emph{IEEE Transactions on Industrial Electronics}, vol. 71, no. 9, pp. 11019--11027, 2024.
\bibitem{ref5} A. ten Pas, M. Gualtieri, K. Saenko and R. Platt, ``Grasp Pose Detection in Point Clouds,'' \emph{International Journal of Robotics Research}, 2017.
\bibitem{ref6} J. Mahler \emph{et al.}, ``Dex-Net 2.0: Deep Learning to Plan Robust Grasps with Synthetic Point Clouds and Analytic Grasp Metrics,'' \emph{Robotics: Science and Systems}, 2017.
\bibitem{ref7} D. Morrison, P. Corke and J. Leitner, ``Closing the Loop for Robotic Grasping: A Real-time, Generative Grasp Synthesis Approach,'' \emph{Robotics: Science and Systems}, 2018.
\bibitem{ref8} A. Mousavian, C. Eppner and D. Fox, ``6-DOF GraspNet: Variational Grasp Generation for Object Manipulation,'' \emph{International Conference on Computer Vision}, 2019.
\bibitem{ref9} H. Liang \emph{et al.}, ``PointNetGPD: Detecting Grasp Configurations from Point Sets,'' \emph{International Conference on Robotics and Automation}, 2019.
\bibitem{ref10} J. Mahler \emph{et al.}, ``Learning Ambidextrous Robot Grasping Policies,'' \emph{Science Robotics}, 2019.
\bibitem{ref11} H.-S. Fang, C. Wang, M. Gou and C. Lu, ``GraspNet-1Billion: A Large-Scale Benchmark for General Object Grasping,'' \emph{Computer Vision and Pattern Recognition}, 2020.
\bibitem{ref12} M. Sundermeyer, A. Mousavian, R. Triebel and D. Fox, ``Contact-GraspNet: Efficient 6-DoF Grasp Generation in Cluttered Scenes,'' \emph{International Conference on Robotics and Automation}, 2021.
\bibitem{ref13} M. Breyer, J. J. Chung, L. Ott, R. Siegwart and J. Nieto, ``Volumetric Grasping Network: Real-time 6 DOF Grasp Detection in Clutter,'' \emph{Conference on Robot Learning}, 2020.
\bibitem{ref14} H.-S. Fang \emph{et al.}, ``AnyGrasp: Robust and Efficient Grasp Perception in Spatial and Temporal Domains,'' \emph{IEEE Transactions on Robotics}, 2023.
\bibitem{ref15} M. Gilles \emph{et al.}, ``MetaGraspNetV2: All-in-One Dataset Enabling Fast and Reliable Robotic Bin Picking via Object Relationship Reasoning and Dexterous Grasping,'' \emph{IEEE Transactions on Automation Science and Engineering}, vol. 21, no. 3, pp. 2302--2320, 2024.
\bibitem{ref16} L. Li, A. Cherouat, H. Snoussi and T. Wang, ``Grasping With Occlusion-Aware Ally Method in Complex Scenes,'' \emph{IEEE Transactions on Automation Science and Engineering}, vol. 22, pp. 5944--5954, 2025.
\bibitem{ref17} X. Zhong \emph{et al.}, ``RIGNet: Robot Intention Grasp for Dense Stacked Targets With Multi-Task Siamese Schema Through RoIs Learning,'' \emph{IEEE Transactions on Automation Science and Engineering}, vol. 22, pp. 10354--10367, 2025.
\bibitem{ref18} Q. Zhao, M. Zheng, Z. Li, S. Huang and W. Shi, ``Robot Dexterous Grasping in Cluttered Scenes Based on Single-View Point Cloud,'' \emph{IEEE Transactions on Automation Science and Engineering}, vol. 22, pp. 17160--17169, 2025.
\bibitem{ref19} G.-H. Xu, Y.-L. Wei, D. Zheng, X.-M. Wu and W.-S. Zheng, ``Dexterous Grasp Transformer,'' \emph{Computer Vision and Pattern Recognition}, pp. 17933--17942, 2024.
\bibitem{ref20} R. Wang \emph{et al.}, ``DexGraspNet: A Large-Scale Robotic Dexterous Grasp Dataset for General Objects Based on Simulation,'' \emph{International Conference on Robotics and Automation}, pp. 11359--11366, 2023.
\bibitem{ref21} P. Li \emph{et al.}, ``GenDexGrasp: Generalizable Dexterous Grasping,'' \emph{International Conference on Robotics and Automation}, 2023.
\bibitem{ref22} Y. Xu \emph{et al.}, ``UniDexGrasp: Universal Robotic Dexterous Grasping via Learning Diverse Proposal Generation and Goal-Conditioned Policy,'' \emph{Computer Vision and Pattern Recognition}, 2023.
\bibitem{ref23} J. Zhang \emph{et al.}, ``DexGraspNet 2.0: Learning Generative Dexterous Grasping in Large-scale Synthetic Cluttered Scenes,'' \emph{Proceedings of the 8th Conference on Robot Learning}, vol. 270, pp. 5106--5133, 2025.
\bibitem{ref24} Z. Weng, H. Lu, D. Kragic and J. Lundell, ``Generating Dexterous Grasps With Diffusion Models,'' \emph{IEEE Robotics and Automation Letters}, 2024.
\bibitem{ref25} Y. Zhong, Q. Jiang, J. Yu and Y. Ma, ``DexGrasp Anything: Towards Universal Robotic Dexterous Grasping with Physics Awareness,'' \emph{Computer Vision and Pattern Recognition}, pp. 22584--22594, 2025.
\bibitem{ref26} J. Jian \emph{et al.}, ``G-DexGrasp: Generalizable Dexterous Grasping Synthesis Via Part-Aware Prior Retrieval and Prior-Assisted Generation,'' \emph{International Conference on Computer Vision}, 2025.
\bibitem{ref27} Z. Chen \emph{et al.}, ``ClutterDexGrasp: A Sim-to-Real System for General Dexterous Grasping in Cluttered Scenes,'' \emph{arXiv preprint arXiv:2506.14317}, 2025.
\bibitem{ref28} P. Grady \emph{et al.}, ``ContactOpt: Optimizing Contact to Improve Grasps,'' \emph{Computer Vision and Pattern Recognition}, 2021.
\bibitem{ref29} H. Jiang \emph{et al.}, ``Hand-Object Contact Consistency Reasoning for Human Grasps Generation,'' \emph{International Conference on Computer Vision}, 2021.
\bibitem{ref30} S. Liu \emph{et al.}, ``ContactGen: Generative Contact Modeling for Grasp Generation,'' \emph{International Conference on Computer Vision}, 2023.
\bibitem{ref31} J. Zhang \emph{et al.}, ``CADGrasp: Learning Contact and Collision Aware General Dexterous Grasping in Cluttered Scenes,'' \emph{arXiv preprint arXiv:2601.15039}, 2026.
\bibitem{ref32} V. Makoviychuk \emph{et al.}, ``Isaac Gym: High Performance GPU-Based Physics Simulation for Robot Learning,'' \emph{Neural Information Processing Systems Datasets and Benchmarks}, 2021.
\bibitem{ref33} Y. Li, W. Wei, D. Li, P. Wang, W. Li and J. Zhong, ``HGC-Net: Deep Anthropomorphic Hand Grasping in Clutter,'' \emph{International Conference on Robotics and Automation}, pp. 714--720, 2022.
\bibitem{ref34} Z. Q. Chen \emph{et al.}, ``Learning Robust Real-World Dexterous Grasping Policies via Implicit Shape Augmentation,'' \emph{Conference on Robot Learning}, 2022.
\bibitem{ref35} A. Kirillov \emph{et al.}, ``Segment Anything,'' \emph{International Conference on Computer Vision}, 2023.

\end{thebibliography}

% Leave unused space at the column bottom, not between author biographies.
\raggedbottom

\begin{IEEEbiography}[{\includegraphics[width=0.78in,height=0.98in,clip,keepaspectratio]{au_hujiarui.png}}]{Jiarui Hu}
received the B.S. degree in mechanical engineering from Shandong University, Jinan, China, in 2021. He is currently working toward the Ph.D. degree in control science and engineering at Tongji University, Shanghai, China. His research interests include tactile sensing and multimodal learning.
\end{IEEEbiography}

\vspace{6pt}
\begin{IEEEbiography}[{\includegraphics[width=0.78in,height=0.98in,clip,keepaspectratio]{au_lixingyu.png}}]{Xingyu Li}
received the M.S. degree in control science and engineering from Zhengzhou University, Zhengzhou, China, in 2024. He is currently working toward the Ph.D. degree in intelligent science and technology with the State Key Laboratory of Autonomous Intelligent Unmanned Systems, Tongji University, Shanghai, China.

His research interests include tactile perception and robot task planning based on large models.
\end{IEEEbiography}

\begin{IEEEbiography}[{\includegraphics[width=0.78in,height=0.98in,clip,keepaspectratio]{au_jiangrong.png}}]{Rong Jiang}
received the B.E. degree in control technology and instrumentation from the University of Electronic Science and Technology of China (UESTC), Chengdu, China, in 2015, the M.S. degree in instrument science and technology from Southeast University, Nanjing, China, in 2018, and the Ph.D. degree in control science and engineering from Tongji University, Shanghai, China, in 2023. He is currently a Postdoctoral Researcher with Tongji University. His research interest is robot skill learning.
\end{IEEEbiography}

\begin{IEEEbiography}[{\includegraphics[width=0.78in,height=0.98in,clip,keepaspectratio]{au_wangzhipeng.png}}]{Zhipeng Wang~(Member, IEEE)}
received the M.S. degree from Zhejiang University, Hangzhou, China, in 2011, and the Ph.D. degree from Tongji University, Shanghai, China, in 2015. From 2015 to 2018, he held a postdoctoral research appointment with the College of Mechanical and Energy Engineering, Tongji University. He is currently an Associate Professor with the Department of Control Science and Engineering, College of Electronics and Information Engineering, Tongji University. His current research interests include autonomous system perception and embodied dexterous manipulation.
\end{IEEEbiography}

\begin{IEEEbiography}[{\includegraphics[width=0.78in,height=0.98in,clip,keepaspectratio]{au_zhouyanmin.png}}]{Yanmin Zhou~(Member, IEEE)}
received the M.S. degree in control theory and control engineering from Tongji University, Shanghai, China, in 2011, and the Ph.D. degree in biomechanics from the University of Cambridge, Cambridge, U.K., in 2015. She is currently an Associate Professor with the Department of Control Science and Engineering, College of Electronics and Information Engineering, Tongji University. Her current research interests include bionics and the design and control of intelligent robots.
\end{IEEEbiography}

\begin{IEEEbiography}[{\includegraphics[width=0.78in,height=0.98in,clip,keepaspectratio]{au_hebin.png}}]{Bin He~(Senior Member, IEEE)}
received the B.S. degree in engineering machinery from Jilin University, Changchun, China, in 1996, and the Ph.D. degree in mechanical and electronic control engineering from Zhejiang University, Hangzhou, China, in 2001. He held a postdoctoral research appointment with the State Key Laboratory of Fluid Power Transmission and Control, Zhejiang University, from 2001 to 2003. He is currently a Professor with the Department of Control Science and Engineering, College of Electronics and Information Engineering, Tongji University, Shanghai, China. His current research interests include intelligent robot control, biomimetic microrobots, and wireless networks.
\end{IEEEbiography}


\end{document}
